% Part: computability
% Chapter: computability-theory
% Section: s-m-n

\documentclass[../../../include/open-logic-section]{subfiles}

\begin{document}

\olfileid{cmp}{thy}{smn}
\olsection{قضیهٔ $s$-$m$-$n$}

\begin{explain}
قضیهٔ بعدی، به دلیلی که به‌زودی روشن می‌شود، «قضیهٔ~$s$-$m$-$n$»
نامیده می‌شود. بخش دشوار، فهم دقیق آن چیزی است که قضیه می‌گوید؛ همین
که گزاره را بفهمید، نسبتاً بدیهی به نظر خواهد رسید.
\end{explain}

\begin{thm}
\ollabel{thm:s-m-n}
  برای هر جفت عدد طبیعی~$n$ و~$m$، تابعی بازگشتی اولیه چون~$s^m_n$
  وجود دارد چنان‌که برای همهٔ اعداد طبیعی
  $e$، $a_0$، \dots، $a_{m-1}$، $y_0$، \dots، $y_{n-1}$، داریم
  \[
  \cfind{s^m_n(e, a_0, \dots, a_{m-1})}[n](y_0, \dots, y_{n-1}) \simeq
  \cfind{e}[m+n](a_0, \dots, a_{m-1}, y_0, \dots, y_{n-1}).
\]
\end{thm}

\begin{explain}
سودمند است~$s^m_n$ را عمل‌کننده بر 
\emph{برنامه‌ها} بدانیم. یعنی $s^m_n$ برنامه‌ای چون~$e$ برای تابعی
$(m+n)$موضعی، همراه با ورودی‌های ثابت
$a_0$، \dots، $a_{m-1}$، می‌گیرد و برنامهٔ
$s^m_n(e,a_0,\dots,a_{m-1})$ را برای تابع~$n$موضعیِ آرگومان‌های
باقی‌مانده بازمی‌گرداند. 
\iftag{TMs}{اگر~$e$ را وصف یک ماشین تورینگ بدانید، آنگاه
$s^m_n(e,a_0,\dots,a_{m-1})$ ماشین تورینگی است که روی ورودی‌های
$y_0$، \dots، $y_{n-1}$، نخست
$a_0$، \dots، $a_{m-1}$ را به رشتهٔ ورودی می‌افزاید و سپس~$e$ را
اجرا می‌کند. در این صورت هر~$s^m_n$ تنها تابعی بازگشتی اولیه است که کد
ماشین تورینگ مناسب را می‌یابد.}{}
\end{explain}

\end{document}

